\begin{figure}[]
    \centering
    % \includegraphics[width=\linewidth]
    % % {images/rdbms_qiskit_best_method.png}
    % {images/162_rdbms_qiskit_best_method2}
    \includegraphics[width=0.7\linewidth]{images/figure7_updated_rdbms_qiskit_winners_new.jpg}
    \caption{\revisionB{Execution time and memory usage comparison between RDBMSs (SQLite, DuckDB, PostgreSQL, and 
Umbra) and Qiskit Aer (state vector, stabilizer, matrix product state, density matrix) over 162 circuits generated by \sys}.
    % Histogram showing percentage of circuits from chosen set that RDBMS methods [sqlite, ducksql, psql from the implementation of Hai et al. (2025) \cite{hai2025quantum}] and Qiskit offered simulator methods [matrix\_product\_state, extended\_stabilizer, density\_matrix] are the best method to simulate based on execution time and memory usage.
    }
    \label{fig:rdbmsqiskitcomparison}
\end{figure}

% \begin{table*}[t]
%   \centering
%   \caption{Routing performance on the 20\% held-out test set. Best linear and tree-based models are bolded.}
%   \label{tab:rdbms_router_combined}
%   \vspace{4pt}
%   % Time-optimal routing
%   \begin{subtable}[t]{0.48\textwidth}
%     \centering
%     \caption{Time-optimal routing}
%     \label{tab:rdbms_time_router} % old label preserved
%     \vspace{2pt}
%     \begin{tabular}{lcccc}
%       \toprule
%       Model & Accuracy & Precision & Recall & F1 \\
%       \midrule
%       Logistic Regression & 0.891 & 0.791 & 0.900 & 0.828 \\
%       \textbf{Linear SVM} & \textbf{0.902} & \textbf{0.806} & \textbf{0.909} & \textbf{0.843} \\
%       \midrule
%       Decision Tree       & 0.934 & 0.876 & 0.876 & 0.876 \\
%       Random Forest       & 0.954 & 0.929 & 0.893 & 0.910 \\
%       \textbf{XGBoost}    & \textbf{0.954} & \textbf{0.909} & \textbf{0.917} & \textbf{0.913} \\
%       \bottomrule
%     \end{tabular}
%   \end{subtable}
%   \hfill
%   % Memory-optimal routing
%   \begin{subtable}[t]{0.48\textwidth}
%     \centering
%     \caption{Memory-optimal routing}
%     \label{tab:rdbms_memory_router} % old label preserved
%     \vspace{2pt}
%     \begin{tabular}{lcccc}
%       \toprule
%       Model & Accuracy & Precision & Recall & F1 \\
%       \midrule
%       Logistic Regression & 0.933 & 0.933 & 0.933 & 0.933 \\
%       \textbf{Linear SVM} & \textbf{0.935} & \textbf{0.935} & \textbf{0.935} & \textbf{0.935} \\
%       \midrule
%       Decision Tree       & 0.956 & 0.956 & 0.956 & 0.956 \\
%       \textbf{Random Forest} & \textbf{0.976} & \textbf{0.976} & \textbf{0.975} & \textbf{0.976} \\
%       XGBoost             & 0.975 & 0.975 & 0.975 & 0.975 \\
%       \bottomrule
%     \end{tabular}
%   \end{subtable}
  
% \end{table*}

% \begin{table*}[t]
%   \centering
%   \caption{Routing performance on the 20\% held-out test set. Best linear and tree-based models are bolded.}
%   \label{tab:rdbms_router_combined}

%   \begin{subtable}[t]{0.48\linewidth}
%     \centering
%     \caption{Time-optimal routing}
%     \label{tab:rdbms_time_router}
%     \begin{tabular}{lcccc}
%       \toprule
%       Model & Accuracy & Precision & Recall & F1 \\
%       \midrule
%       Logistic Regression & 0.887 & 0.775 & 0.908 & 0.817 \\
%       \textbf{Linear SVM} & \textbf{0.908} & \textbf{0.803} & \textbf{0.918} & \textbf{0.843} \\
%       \midrule
%       Decision Tree & 0.940 & 0.881 & 0.872 & 0.877 \\
%       \textbf{Random Forest} & \textbf{0.956} & \textbf{0.919} & \textbf{0.897} & \textbf{0.907} \\
%       XGBoost & 0.953 & 0.892 & 0.923 & 0.907 \\
%       \bottomrule
%     \end{tabular}
%   \end{subtable}
%   \hfill
%   \begin{subtable}[t]{0.48\linewidth}
%     \centering
%     \caption{Memory-optimal routing}
%     \label{tab:rdbms_memory_router}
%     \begin{tabular}{lcccc}
%       \toprule
%       Model & Accuracy & Precision & Recall & F1 \\
%       \midrule
%       Logistic Regression & 0.935 & 0.932 & 0.934 & 0.933 \\
%       \textbf{Linear SVM} & \textbf{0.940} & \textbf{0.937} & \textbf{0.938} & \textbf{0.938} \\
%       \midrule
%       Decision Tree & 0.959 & 0.960 & 0.956 & 0.958 \\
%       Random Forest & 0.973 & 0.975 & 0.970 & 0.972 \\
%       \textbf{XGBoost} & \textbf{0.976} & \textbf{0.978} & \textbf{0.973} & \textbf{0.975} \\
%       \bottomrule
%     \end{tabular}
%   \end{subtable}

% \end{table*}


\begin{table*}[t]
  \centering
  \caption{Routing performance on the 20\% held-out test set. Models are evaluated using accuracy and normalized confusion-matrix entries. The columns represent true label $\rightarrow$ prediction label. 
  % Best linear and tree-based models are bolded.
  }
  \label{tab:rdbms_router_combined}
 % \vspace{-0.4cm}
  % \begin{subtable}[t]{0.48\linewidth}
  %   \centering
  %   \caption{Time-optimal routing}
  %   \label{tab:rdbms_time_router}
  %   \resizebox{\linewidth}{!}{
  %   \begin{tabular}{lccccc}
  %     \toprule
  %     Model & Accuracy & rdbms$\rightarrow$rdbms & qiskit$\rightarrow$qiskit & rdbms$\rightarrow$qiskit & qiskit$\rightarrow$rdbms \\
  %     \midrule
  %     Logistic Regression & 0.887 & \textbf{0.94} & 0.88 & 0.06 & \textbf{0.12}  \\
  %     \textbf{Linear SVM} & \textbf{0.908} & 0.93 & \textbf{0.90} & \textbf{0.07} & 0.10\\
  %     \midrule
  %     Decision Tree & 0.940 & 0.78 & 0.97 & \textbf{0.22} & 0.03 \\
  %     Random Forest & \textbf{0.956} & 0.81 & 0.98 & 0.19 & 0.02 \\
  %     \textbf{XGBoost} & 0.953 & \textbf{0.88} & \textbf{0.96} & 0.12 & \textbf{0.04} \\
  %     \bottomrule
  %   \end{tabular}}
  % \end{subtable}
  % \hfill
  % \begin{subtable}[t]{0.48\linewidth}
  %   \centering
  %   \caption{Memory-optimal routing}
  %   \label{tab:rdbms_memory_router}
  %   \resizebox{\linewidth}{!}{
  %   \begin{tabular}{lccccc}
  %     \toprule
  %     Model & Accuracy & rdbms$\rightarrow$rdbms & qiskit$\rightarrow$qiskit & rdbms$\rightarrow$qiskit & qiskit$\rightarrow$rdbms \\
  %     \midrule
  %     Logistic Regression & 0.935 & 0.93 & 0.94 & 0.07 & 0.06 \\
  %     \textbf{Linear SVM} & \textbf{0.940} & \textbf{0.93} & \textbf{0.95} & \textbf{0.07} & \textbf{0.05} \\
  %     \midrule
  %     Decision Tree & 0.962 & 0.94 & 0.98 & 0.06 & 0.02 \\
  %     \textbf{Random Forest} & \textbf{0.973} & \textbf{0.95} & \textbf{0.99} & \textbf{0.05} & \textbf{0.01} \\
  %     \textbf{XGBoost} & \textbf{0.974} & \textbf{0.95} & \textbf{0.99} & \textbf{0.05} & \textbf{0.01} \\
  %     \bottomrule
  %   \end{tabular}}
  % \end{subtable}
 \begin{subtable}[t]{0.48\linewidth}
    \centering
    \caption{Time-optimal routing}
    \label{tab:rdbms_time_router}
    \resizebox{\linewidth}{!}{
    \begin{tabular}{lccccc}
      \toprule
      Model & Accuracy & rdbms$\rightarrow$rdbms & qiskit$\rightarrow$qiskit & rdbms$\rightarrow$qiskit & qiskit$\rightarrow$rdbms \\
      \midrule
      Logistic Regression & 0.887 & 0.94 & 0.88 & 0.06 & 0.12  \\
      Linear SVM & 0.908 & 0.93 & 0.90 & 0.07 & 0.10\\
      \midrule
      Decision Tree & 0.940 & 0.78 & 0.97 & 0.22 & 0.03 \\
      Random Forest & 0.956 & 0.81 & 0.98 & 0.19 & 0.02 \\
      XGBoost & 0.953 & 0.88 & 0.96 & 0.12 & 0.04 \\
      \bottomrule
    \end{tabular}}
  \end{subtable}
  \hfill
  \begin{subtable}[t]{0.48\linewidth}
    \centering
    \caption{Memory-optimal routing}
    \label{tab:rdbms_memory_router}
    \resizebox{\linewidth}{!}{
    \begin{tabular}{lccccc}
      \toprule
      Model & Accuracy & rdbms$\rightarrow$rdbms & qiskit$\rightarrow$qiskit & rdbms$\rightarrow$qiskit & qiskit$\rightarrow$rdbms \\
      \midrule
      Logistic Regression & 0.935 & 0.93 & 0.94 & 0.07 & 0.06 \\
      Linear SVM & 0.940 & 0.93 & 0.95 & 0.07 & 0.05 \\
      \midrule
      Decision Tree & 0.962 & 0.94 & 0.98 & 0.06 & 0.02 \\
      Random Forest & 0.973 & 0.95 & 0.99 & 0.05 & 0.01 \\
      XGBoost & 0.974 & 0.95 & 0.99 & 0.05 & 0.01 \\
      \bottomrule
    \end{tabular}}
  \end{subtable}

\end{table*}




\subsection{Use case 1: RDBMS as Simulation Backend}
\label{ssec:simu}
\emph{Q1: Is an RDBMS an efficient simulation engine?}

To answer this question, we generated 162 circuits\footnote{\url{https://github.com/InfiniData-Lab/InferQ/blob/main/analysis/training_data/rdbms_all_methods_training_data.parquet}} using \sys. For each circuit, we
measured wall-clock runtime and peak memory when simulating it with (i) 
RDBMS engines (PostgreSQL, SQLite, DuckDB, and \revisionB{Umbra}) executing the SQL workload emitted by
\sys, and (ii) Qiskit Aer using four built-in simulation methods\footnote{Qiskit Aer Simulation Methods: \url{https://qiskit.github.io/qiskit-aer/tutorials/1_aersimulator.html##Simulation-Method-Option}}
(\feat{statevector}, \feat{extended_stabilizer}, \feat{matrix_product_state}, and
\feat{density_matrix}).
% \aad{plz add here with state vector, and link to each of them}


Figure~\ref{fig:rdbmsqiskitcomparison} reports, for each metric, the fraction of circuits
for which a backend achieves the best result (minimum execution time or minimum peak
memory) among all evaluated options. As expected, Qiskit Aer attains the best execution
time on most circuits. 
% More interestingly, even in this untuned setting, 
SQLite
achieves the lowest peak memory on 66.7\% of the circuits (108/162). This suggests that
RDBMS engines can be competitive simulation backends, particularly for \emph{memory}. This result 
motivates  further work on SQL workload optimization. 
Moreover, it is interesting to explore learning-based routing in simulation, i.e., predicting when an RDBMS engine is likely to outperform conventional simulators. This leads to our next question.

 
 

% \begin{table}[t]
%   \centering
%   \caption{Runtime (time)-based routing performance on the held-out test set.}
%   \label{tab:rdbms_time_router}
%   \vspace{2pt}
%   \begin{tabular}{lcccc}
%     \toprule
%     Model & Accuracy & Precision & Recall & F1 \\
%     \midrule
%     Logistic Regression & 0.876 & 0.542 & 0.892 & 0.675 \\
%     Linear SVM          & 0.894 & 0.587 & 0.901 & 0.710 \\
%     Random Forest       & 0.942 & 0.820 & 0.761 & 0.790 \\
%     XGBoost             & 0.941 & 0.779 & 0.824 & 0.801 \\
%     \bottomrule
%   \end{tabular}
% \end{table}

% \begin{figure}[t]
%     \centering
%     % \includegraphics[width=\linewidth]{images/timeoptallrdbms.png}
%     \caption{Confusion plots for Linear SVM and XGBoost models trained for time-optimal routing between the SQLite RDBMS backend and Qiskit Aer using InferQ numeric features.}
%     \label{fig:rdbmsqiskitimportancemem}
% \end{figure}
% \begin{table}[t]
%   \centering
%   \caption{Memory-based routing performance on the held-out test set.}
%   \label{tab:rdbms_memory_router}
%   \vspace{2pt}
%   \begin{tabular}{lcccc}
%     \toprule
%     Model & Accuracy & Precision & Recall & F1 \\
%     \midrule
%     Logistic Regression & 0.942 & 0.925 & 0.938 & 0.931 \\
%     Linear SVM          & 0.943 & 0.929 & 0.935 & 0.932 \\
%     Random Forest       & 0.979 & 0.995 & 0.953 & 0.974 \\
%     XGBoost             & 0.979 & 0.990 & 0.958 & 0.974 \\
%     \bottomrule
%   \end{tabular}
% \end{table}
% \begin{figure}[h]
%     \centering
%     % \includegraphics[width=\linewidth]{images/memoptallrdbms.png}
%     \caption{Confusion plots for Linear SVM and Random Forest models trained for memory-optimal routing between the SQLite RDBMS backend and Qiskit Aer using InferQ numeric features.}
%     \label{fig:rdbmsqiskitimportancemem}
% \end{figure}
 % \begin{figure}[t]
 %  \centering
 %   \includegraphics[width=\linewidth]{images/qft_26q_cap_sensitivity_silverpies.pdf}
 %    \vspace{-1mm}
 %    \caption{26q cap sensitivity.}
 %    \label{fig:qft-26q-cap-sensitivity}
 %    \end{figure}

 %     \begin{figure}[t]
 %  \centering
 %   \includegraphics[width=\linewidth]{images/qft_23_26_evolution_silverpies.pdf}
 %    \vspace{-1mm}
 %    \caption{Evolution from 23 to 26 qubits.}
 %    \label{fig:qft-26q-cap-sensitivity}
 %    \end{figure}
% \begin{figure*}[t]
%   \centering

%   \newlength{\FigGap}
%   \setlength{\FigGap}{0.03\textwidth}

%   \newlength{\LeftH}
%   \newlength{\RightH}
%   \setlength{\LeftH}{4.0cm}
%   \setlength{\RightH}{4.1cm}
%   \begin{minipage}[t]{0.55\textwidth}
%     \vspace{0pt}
%     \centering
%     \includegraphics[
%       width=\linewidth,
%       height=\RightH,
%       keepaspectratio
%     ]{images/qft_23_26_evolution_silverpies.pdf}
%     \vspace{-3mm}
%     \captionof{subfigure}{\revisionA{Runtime and spill increase from 23 to 26 qubits for QFT circuits.}}
%     \label{fig:qft-23-26-evolution}
%   \end{minipage}
%   \hspace{\FigGap}%
% \begin{minipage}[t]{0.395\textwidth}
%     \vspace{0pt}
%     \centering
%     \includegraphics[
%       width=\linewidth,
%       height=\LeftH,
%       keepaspectratio
%     ]{images/qft_26q_cap_sensitivity_silverpies.pdf}
%     \vspace{-6mm}
%     \captionof{subfigure}{\revisionA{Runtime and spill volume for 26-qubit QFT under different memory limits for QFT circuits.}}
%     \label{fig:qft-26q-cap-sensitivity}
%   \end{minipage}%

%   \vspace{-1mm}

% \end{figure*}

% \begin{figure*}[t]
%   \centering

%   \newlength{\FigGap}
%   \setlength{\FigGap}{0.01\textwidth}

%   \newlength{\LeftH}
%   \newlength{\RightH}
%   \setlength{\LeftH}{4.0cm}
%   \setlength{\RightH}{4.0cm}
%   \begin{minipage}[t]{0.49\textwidth}
%     \vspace{0pt}
%     \centering
%     \includegraphics[
%       width=\linewidth,
%       height=\RightH,
%       keepaspectratio
%     ]{images/fig2_spill_vs_qubits.pdf}
%     \vspace{-3mm}
%     \captionof{subfigure}{\revisionA{Spill volume across circuit computation plotted against the number of qubits of the circuit samples (\sys circuits). Spill is plotted on log scale.}}
%     \label{fig:spill_vs_qubits}
%   \end{minipage}
%   \hspace{\FigGap}%
% \begin{minipage}[t]{0.49\textwidth}
%     \vspace{0pt}
%     \centering
%     \includegraphics[
%       width=\linewidth,
%       height=\LeftH,
%       keepaspectratio
%     ]{images/fig7_cte_vs_density.pdf}
%     \vspace{-6mm}
%     \captionof{subfigure}{\revisionA{CTE sizes total and largest across circuit computation plotted against the sparsity (density in InferQ) of the circuit samples. CTE sizes are plotted on log scale.}}
%     \label{fig:cte_vs_density}
%   \end{minipage}%

%   \vspace{-1mm}

% \end{figure*}
% Fig. 8


%-----------------------------------------------------------------------------
\begin{figure}[t]
  \centering
  \includegraphics[width=0.7\columnwidth]{images/expanderfig3_spill_vs_cte.pdf}
  \caption{\revisionA{Out-of-core spill versus intermediate relation sizes.}}
  \label{fig:ooc-expander-spill-cte}
\end{figure}

% Fig. 9
\begin{figure}[t]
  \centering
  \includegraphics[width=0.65\columnwidth]{images/fig4a_walltime_vs_qubits.pdf}
  \caption{\revisionA{Runtime versus \#qubits.}}
  \label{fig:ooc-expander-runtime-qubits}
\end{figure}

% Fig. 10
\begin{figure}[t]
  \centering
  \includegraphics[width=0.65\columnwidth]{images/fig2_spill_vs_qubits.pdf}
    % \vspace{-0.4cm}
  \caption{\revisionA{Spill versus \#qubits on sampled InferQ circuits.}}
  \label{fig:ooc-inferq-spill-qubits}
   \vspace{-0.4cm}
\end{figure}

% Fig. 11
\begin{figure}[t]
  \centering
  \includegraphics[width=0.7\columnwidth]{images/fig7_cte_vs_density.pdf}
    \vspace{-0.4cm}
  \caption{\revisionA{Intermediate CTE sizes versus output density on sampled InferQ circuits.}}
  \label{fig:ooc-inferq-cte-density}
    \vspace{-0.4cm}
\end{figure}
 

% Optional: enforce consistent panel heights for a clean SIGMOD 2-col look.
% \newcommand{\oocPanelH}{3.2cm} % adjust (e.g., 3.0--3.5cm) if needed

% % ---------------- Figure 8 (Expander) ----------------
% \begin{figure*}[t]
%   \centering
%   \begin{subfigure}[t]{0.49\textwidth}
%     \centering
%     \includegraphics[width=\linewidth,height=\oocPanelH,keepaspectratio]{images/expanderfig3_spill_vs_cte.pdf}
%     \caption{\protect\revisionA{Spill vs intermediate sizes.}}
%     \label{fig:ooc-expander-spill-cte}
%   \end{subfigure}\hfill
%   \begin{subfigure}[t]{0.49\textwidth}
%     \centering
%     \includegraphics[width=\linewidth,height=\oocPanelH,keepaspectratio]{images/fig4a_walltime_vs_qubits.pdf}
%     \caption{\protect\revisionA{Runtime vs \#qubits.}}
%     \label{fig:ooc-expander-runtime-qubits}
%   \end{subfigure}
%   \caption{\protect\revisionA{Out-of-core results on Expander circuits under a 16~GB memory limit.}}
%   \label{fig:ooc-expander}
% \end{figure*}

% % ---------------- Figure 9 (InferQ samples) ----------------
% \begin{figure*}[t]
%   \centering
%   \begin{subfigure}[t]{0.49\textwidth}
%     \centering
%     \includegraphics[width=\linewidth,height=\oocPanelH,keepaspectratio]{images/fig2_spill_vs_qubits.pdf}
%     \caption{\protect\revisionA{Spill vs \#qubits.}}
%     \label{fig:ooc-inferq-spill-qubits}
%   \end{subfigure}\hfill
%   \begin{subfigure}[t]{0.49\textwidth}
%     \centering
%     \includegraphics[width=\linewidth,height=\oocPanelH,keepaspectratio]{images/fig7_cte_vs_density.pdf}
%     \caption{\protect\revisionA{CTE sizes vs density.}}
%     \label{fig:ooc-inferq-cte-density}
%   \end{subfigure}
%   \caption{\protect\revisionA{Out-of-core behavior across output-density bands on sampled InferQ circuits (16~GB).}}
%   \label{fig:ooc-inferq}
% \end{figure*}


  \vspace{0.2cm}
\noindent\emph{Q2: How to decide whether to choose an RDBMS as a simulation engine?} %How should the choice of an RDBMS as a simulation engine be decided? 
  \vspace{0.2cm}

We model this choice as a \emph{binary selection} problem: given a circuit, predict
whether an RDBMS engine achieves lower runtime and/or lower peak memory than a
conventional simulator. Prior database-oriented studies of SQL-based simulation report
strong results on a small set of highly structured circuits (e.g., sparse state-preparation
workloads)~\cite{hai2025quantum,einsteinSQL2023}. Using \sys, we move beyond these
constrained cases and quantify when an RDBMS is beneficial over a substantially broader
set of 
% \emph{general, compositional circuits}.
general circuits.
%
 We have created a training dataset of \num{7705} circuits\footnote{Training data (Parquet) available at \url{https://github.com/InfiniData-Lab/InferQ/blob/main/analysis/training_data/rdbms_training_data.parquet}. The results reported in Figure~\ref{fig:aer-vs-sqlite-wins} also used this dataset.} and train five standard 
% \rh{linear and tree} 
models as selectors: logistic regression,
linear SVM,  Decision Tree, random forest (RF), and XGBoost. 
For each circuit, we have logged runtime and peak
memory for Qiskit Aer and the RDBMS backends. This yields ground-truth labels
(\emph{RDBMS wins} vs.\ \emph{Qiskit wins}) for both objectives.

Table~\ref{tab:rdbms_time_router} shows that
these models achieve high accuracy using the features extracted in
Section~\ref{ssec:features}, despite class imbalance (since SQLite is faster only for a
subset of circuits). 
% XGBoost performs best overall (95.3\% accuracy, lower confusion than RF), consistent with its
% effectiveness for learned cost models~\cite{adams2019learning,chen2018tvm}. 
Random Forest achieves the highest runtime-routing accuracy (95.6\%), while XGBoost follows closely (95.3\%) and better identifies RDBMS-win cases.
The best linear model, linear SVM, too achieves 90.8\% accuracy with very low off diagonal confusion entries on the test set, suggesting that much of the
time-selection boundary is close to linear in our feature space. Given the strong performance of these lightweight models, we did not evaluate more
complex models that require substantially higher training and tuning cost.
% Overall, \sys features make the RDBMS selection decision predictable with simple,
% practical models. 

For memory efficiency
(Table~\ref{tab:rdbms_memory_router}), performance is similarly strong: XGBoost achieves the highest accuracy (97.4\%), while random forest follows closely (97.3\%). Linear SVM also remains  highly competitive (94.0\%
accuracy). Overall, \sys features in Section~\ref{ssec:features} make the RDBMS selection decision
predictable with simple, practical models.

 


 
\revisionA{
\para{Routing overhead}
The learned selector is trained offline on labeled workloads generated by \sys%InferQ
and deployed only for inference. 
% Thus, it is outside the inner simulation loop: at routing time, the added cost is pre-simulation feature extraction plus one classifier evaluation. For strict online routing, the selector can use static, graph, and SQL features, whose ablation performance remains high (Table~\ref{tab:feature-ablation}). 
For XGBoost, which provides the best accuracy--efficiency trade-off among the models in Tables~\ref{tab:rdbms_time_router} and~\ref{tab:rdbms_memory_router}, inference takes only 13.9\,ms for runtime routing and 13.5\,ms for memory routing per batch.  Appendix~C %\ref{sec:traintime}  
 of our technical report \cite{inferqtech} reports the full training and inference times for all models. Our goal is therefore not to add heavyweight ML components to simulation, but to study a lightweight learned routing policy for backend selection, aligned with recent database work on learned decision-making and cost models~\cite{DBLP:journals/pacmmod/HeinrichLWKB25, DBLP:journals/pacmmod/YangWZDLC23}.}
 

 

\revisionA{The previous experiments compare runtime and peak memory only when all backends finish successfully. For larger circuits, conventional in-memory simulators such as Qiskit Aer can fail once the state representation exceeds available memory. RDBMS engines can instead execute queries out of core by spilling intermediate results to disk. This motivates the following question:}

 


  \vspace{0.2cm}
% \revisionA{\emph{Q3: Can RDBMS-based simulation still complete when the workload no longer fits in memory?
% }}
\noindent\revisionA{\emph{Q3: Can RDBMS-based simulation complete workloads that no longer fit in memory?}}
 % can SQL-driven simulation still complete when the workload no longer fits in memory?
  \vspace{0.2cm}
  
% \revisionA{
% We evaluate this question using QFT circuits, which are dense and quickly produce large intermediate states. We execute SQL workloads for QFT circuits up to 26 qubits on PostgreSQL, DuckDB, and SQLite under 4, 8, and 16\,GB memory limits, using one CPU core. Figure~\ref{fig:qft-23-26-evolution} shows the runtime and spill evolution from 23 to 26 qubits, where spilling becomes increasingly visible. Figure~\ref{fig:qft-26q-cap-sensitivity} summarizes the 26-qubit case, where Qiskit Aer fails with an out-of-memory error, while several RDBMS configurations still complete by spilling to disk.  Full experimental settings and detailed results, including additional \sys-generated circuits, are reported in Appendix~A.
% % Each configuration is run once for warm-up and five times for measurement 
% }

% \revisionA{
% The results show that out-of-core execution is a concrete advantage of RDBMS-backed simulation. For the 26-qubit QFT workload, PostgreSQL completes all memory limits and remains spill-heavy; DuckDB is highly sensitive to the memory limit; and SQLite is slower but robust. These engine-specific profiles show that \sys can expose not only when an RDBMS is faster or more memory-efficient, but also when it can complete simulations that Qiskit Aer cannot.
% }

 

 

\revisionA{We evaluate Q3 using a family of large, sparse circuits with up to 45 qubits under a 16~GB memory limit and one CPU core. Full experimental details are in Appendix~A of   our technical report. Under this memory limit, Qiskit Aer’s exact \texttt{statevector} method fails from 30 qubits onward, while PostgreSQL, DuckDB, and SQLite complete all runs by spilling temporary data to disk. Figure~\ref{fig:ooc-expander-spill-cte} shows that intermediate relations remain MB-scale, but spill reaches GB-scale and correlates more with the total intermediate volume than with the single largest intermediate, indicating that cumulative joins and group-bys lead to externalization. Figure~\ref{fig:ooc-expander-runtime-qubits} reports runtime: PostgreSQL is fastest, DuckDB is second, and SQLite is slower but robust. We have also observed that runtime is strongly correlated with spill volume. 
%  (Appendix~A).
}

\revisionA{To understand out-of-core performance and the characteristics of the simulation workload, we have also sampled 17 InferQ circuits (20--23 qubits) with different output state sparsity. Figures~\ref{fig:ooc-inferq-spill-qubits} and~\ref{fig:ooc-inferq-cte-density} show that sparse outputs usually incur less spill, but output state sparsity alone does not determine spill: several medium-density circuits spill more than dense ones due to intermediate-relation growth and optimizer choices (e.g., join ordering and materialization). Full settings and additional results are reported in \cite{inferqtech}, Appendix~A.}




\medskip
\noindent
\setlength{\fboxsep}{6pt}
\fcolorbox{black!15}{black!4}{%
  \parbox{\dimexpr\linewidth-2\fboxsep-2\fboxrule\relax}{%
\para{Take-away}
\textbf{(1)} Across a wide range of \emph{general, compositional} circuits, RDBMS engines can be
competitive simulation engines, especially in peak memory, and in some cases also in runtime.
\emph{This enables follow-up work on optimizing the emitted SQL to further improve RDBMS performance.}\\[2pt]
\textbf{(2)} \sys enables data-centric routing at scale: it generates diverse circuits, labels each circuit
with observed runtime and peak memory, and exposes features that allow simple models to decide when an
RDBMS engine should be used.
\emph{This provides a foundation for learning-based simulator selection and optimization, trained on \sys-generated circuits.}

% \revisionA{\textbf{(3)} For large circuits that exceed memory, RDBMS engines can still complete simulation through out-of-core execution.
% \emph{This highlights a distinct advantage of database-backed simulation and enables follow-up work on spill-aware SQL plans and memory-sensitive query optimization.}}

\revisionA{(3) For large circuits that exceed memory, RDBMS backends can still complete simulation via out-of-core execution. \emph{This highlights a distinct advantage of database-backed simulation and enables follow-up work on spill-aware planning and optimization.}
}}}
\vspace{0.2cm}
%  \para{Take-away}
% (1) Across a wide range of \emph{general, compositional} circuits, RDBMS engines can be
% competitive simulation engines, especially in peak memory, and in some cases also in
% runtime. (2) \sys enables data-centric routing at scale: it generates diverse circuits,
% provides labels from measured performance of runtime and peak memory, and exposes features that allow
% simple models to accurately decide when an RDBMS engine should be used.
% We train a binary router using InferQ numeric circuit and graph-derived features to predict whether the SQLite RDBMS backend outperforms the best Qiskit Aer method in \emph{runtime}, using 7,713 labeled circuits.




% From Figure \ref{fig:rdbmsqiskitcomparison}, it is clear that only SQLite will be considered as the RDBMS system option. Using all \textit{numeric} InferQ features for the circuits we can train a selector for deciding whether an RDBMS (SQLite) system would be more resource efficient for simulating a particular circuit. To do this we use a subset of 7713 InferQ circuits which have simulation resources recorded for RDBMS (SQLite) and at least 1 Qiskit Aer method from the 4 in Figure \ref{fig:rdbmsqiskitcomparison}.


 
% \begin{table*}[t]
%   \centering
%   \caption{Routing performance using only SQL-derived features on the 20\% held-out test set. Best linear and tree-based models for each routing objective are shown in bold. XG boost chosen over Random Forest due to less confusion % \rh{Do we need commented out fig:xgrftimeconfusion to support this, maybe in appendix?}, despite having better F1-score. 
%   }
%   \label{tab:sql_only_routing}
%   \vspace{4pt}
% % Time-optimal routing
%   \begin{subtable}[t]{0.48\textwidth}
%     \centering
%     \caption{Time-optimal routing}
%     \label{tab:sql_only_time} % sublabel preserved
%     \vspace{2pt}
%     \begin{tabular}{lcccc}
%       \toprule
%       Model & Accuracy & Precision & Recall & F1 \\
%       \midrule
%       Logistic Regression & 0.787 & 0.694 & 0.826 & 0.712 \\
%       \textbf{Linear SVM} & \textbf{0.807} & \textbf{0.707} & \textbf{0.834} & \textbf{0.731} \\
%       \midrule
%       Decision Tree       & 0.903 & 0.812 & 0.837 & 0.824 \\
%       Random Forest       & 0.925 & 0.862 & 0.850 & 0.856 \\
%       \textbf{XGBoost} & \textbf{0.915} & \textbf{0.829} & \textbf{0.879} & \textbf{0.851} \\
%       \bottomrule
%     \end{tabular}
%   \end{subtable}
%   \hfill
%   % Memory-optimal routing
%   \begin{subtable}[t]{0.48\textwidth}
%     \centering
%     \caption{Memory-optimal routing}
%     \label{tab:sql_only_memory} % sublabel preserved
%     \vspace{2pt}
%     \begin{tabular}{lcccc}
%       \toprule
%       Model & Accuracy & Precision & Recall & F1 \\
%       \midrule
%       Logistic Regression & 0.756 & 0.758 & 0.757 & 0.756 \\
%       \textbf{Linear SVM} & \textbf{0.761} & \textbf{0.763} & \textbf{0.762} & \textbf{0.761} \\
%       \midrule
%       Decision Tree       & 0.805 & 0.805 & 0.805 & 0.805 \\
%       \textbf{Random Forest} & \textbf{0.844} & \textbf{0.844} & \textbf{0.844} & \textbf{0.844} \\
%       XGBoost             & 0.844 & 0.843 & 0.843 & 0.843 \\
%       \bottomrule
%     \end{tabular}
%   \end{subtable}
  
% \end{table*}


% \begin{table*}[t]
%   \centering
%   \caption{Routing performance using only SQL-derived features on the 20\% held-out test set. Models are evaluated using accuracy and normalized confusion-matrix entries. The columns represent true label $\rightarrow$ prediction label. 
%   Best linear and tree-based models are bolded.
%   }
% \label{tab:sql_only_router_combined}
%   \vspace{-0.2cm}
%   \begin{subtable}[t]{0.48\linewidth}
%     \centering
%     \caption{Time-optimal routing}
%     \label{tab:sql_only_routing}
%     \resizebox{\linewidth}{!}{
%     \begin{tabular}{lccccc}
%       \toprule
%       Model & Accuracy & rdbms$\rightarrow$rdbms & qiskit$\rightarrow$qiskit & rdbms$\rightarrow$qiskit & qiskit$\rightarrow$rdbms \\
%       \midrule
%       Logistic Regression & 0.787 & 0.88 & 0.77 & 0.12 & 0.23 \\
%       \textbf{Linear SVM} & \textbf{0.803} & \textbf{0.88} & \textbf{0.79} & \textbf{0.12} & \textbf{0.21} \\
%       \midrule
%       Decision Tree & 0.915 & 0.76 & 0.94 & 0.24 & 0.06 \\
%       Random Forest & 0.932 & 0.75 & 0.96 & 0.25 & 0.04 \\
%       \textbf{XGBoost} & \textbf{0.924} & \textbf{0.81} & \textbf{0.94} & \textbf{0.19} & \textbf{0.06} \\
%       \bottomrule
%     \end{tabular}}
%   \end{subtable}
%   \hfill
%   \begin{subtable}[t]{0.48\linewidth}
%     \centering
%     \caption{Memory-optimal routing}
%     \label{tab:sql_only_memory}
%     \resizebox{\linewidth}{!}{
%     \begin{tabular}{lccccc}
%       \toprule
%       Model & Accuracy & rdbms$\rightarrow$rdbms & qiskit$\rightarrow$qiskit & rdbms$\rightarrow$qiskit & qiskit$\rightarrow$rdbms \\
%       \midrule
%       Logistic Regression & 0.735 & 0.77 & 0.71 & 0.23 & 0.29 \\
%       \textbf{Linear SVM} & \textbf{0.737} & \textbf{0.77} & \textbf{0.72} & \textbf{0.23} & \textbf{0.28} \\
%       \midrule
%       Decision Tree & 0.836 & 0.81 & 0.85 & 0.19 & 0.15 \\
%       \textbf{Random Forest} & \textbf{0.866} & \textbf{0.82} & \textbf{0.90} & \textbf{0.18} & \textbf{0.10} \\
%       XGBoost & 0.842 & 0.84 & 0.85 & 0.16 & 0.15 \\
%       \bottomrule
%     \end{tabular}}
%   \end{subtable}

% \end{table*}
\begin{table*}[t]
  \centering
  \caption{Routing performance using only SQL-derived features on the 20\% held-out test set. Models are evaluated using accuracy and normalized confusion-matrix entries. The columns represent true label $\rightarrow$ prediction label. 
  % Best linear and tree-based models are bolded.
  }
\label{tab:sql_only_router_combined}
  \vspace{-0.2cm}
  \begin{subtable}[t]{0.48\linewidth}
    \centering
    \caption{Time-optimal routing}
    \label{tab:sql_only_routing}
    \resizebox{\linewidth}{!}{
    \begin{tabular}{lccccc}
      \toprule
      Model & Accuracy & rdbms$\rightarrow$rdbms & qiskit$\rightarrow$qiskit & rdbms$\rightarrow$qiskit & qiskit$\rightarrow$rdbms \\
      \midrule
      Logistic Regression & 0.787 & 0.88 & 0.77 & 0.12 & 0.23 \\
      Linear SVM & 0.803 & 0.88 & 0.79 & 0.12 & 0.21 \\
      \midrule
      Decision Tree & 0.915 & 0.76 & 0.94 & 0.24 & 0.06 \\
      Random Forest & 0.932 & 0.75 & 0.96 & 0.25 & 0.04 \\
      XGBoost & 0.924 & 0.81 & 0.94 & 0.19 & 0.06 \\
      \bottomrule
    \end{tabular}}
  \end{subtable}
  \hfill
  \begin{subtable}[t]{0.48\linewidth}
    \centering
    \caption{Memory-optimal routing}
    \label{tab:sql_only_memory}
    \resizebox{\linewidth}{!}{
    \begin{tabular}{lccccc}
      \toprule
      Model & Accuracy & rdbms$\rightarrow$rdbms & qiskit$\rightarrow$qiskit & rdbms$\rightarrow$qiskit & qiskit$\rightarrow$rdbms \\
      \midrule
      Logistic Regression & 0.735 & 0.77 & 0.71 & 0.23 & 0.29 \\
      Linear SVM & 0.737 & 0.77 & 0.72 & 0.23 & 0.28 \\
      \midrule
      Decision Tree & 0.836 & 0.81 & 0.85 & 0.19 & 0.15 \\
      Random Forest & 0.866 & 0.82 & 0.90 & 0.18& 0.10 \\
      XGBoost & 0.842 & 0.84 & 0.85 & 0.16 & 0.15 \\
      \bottomrule
    \end{tabular}}
  \end{subtable}

\end{table*}


\subsection{Feature Validation}
\label{ssec:feat_val}












% \begin{figure}[h]
%     \centering
%     \includegraphics[width=\linewidth]{images/xgrftimesqlconfusion.png}
%     \caption{ \rh{Is this Confusion plot necessary? Or can we just state it in the caption of table 9 why we bolded XG boost over RF} XG boost having better confusion results over RF for time optimality training with SQL features only due to class imbalance.}
%     \label{fig:xgrftimeconfusion}
% \end{figure}

Next, we want to go deeper and understand:

\noindent\emph{Q4: Which properties determine whether an RDBMS engine is an efficient simulation engine?}

Our goal is to validate that the four feature groups introduced in Section~\ref{ssec:features}
(\emph{static, graph, dynamic, SQL}) capture the key signals behind the routing outcome in
Section~\ref{ssec:simu}. 
% We therefore analyze feature importance for the best models:
% XGBoost for \emph{runtime} and Random Forest for \emph{peak memory}. 
We therefore analyze feature importance for two high-performing tree-based models: XGBoost for \emph{runtime} and Random Forest for \emph{peak memory}.
We additionally report a
Linear SVM to expose the \emph{direction} of influence via signed coefficients. \revisionB{We also include an ablation study.}




\subsubsection{SQL-only features.}
\label{ssec:sqlonly}
We first restrict to SQL-derived features extracted from output SQL queries representing the generated circuits
(Section~\ref{ssec:sql}, Table~\ref{tab:sql_complexity}).
Table~\ref{tab:sql_only_routing} shows that these SQL-only features already yield strong routing
accuracy, especially for time-optimal routing (XGBoost exceeds 92\% accuracy on the held-out
test set). This indicates that the engine choice is strongly correlated with \emph{query shape},
without requiring circuit-specific features. \revisionM{We have included more results on dominant SQL signals in Appendix~L %\ref{ssec:sqlf} 
of \cite{inferqtech}}.

\begin{figure}[t]
    \centering
    \includegraphics[width=0.8\linewidth]{images/feature_importance_all_features_time_rdbms_vs_qiskit.jpg}
    \caption{Importance of top 10 InferQ numeric features for SVM and XGBoost trained for optimizing execution time.}
    \label{fig:timeoptallrdbmsimp}
\end{figure}
 \begin{figure}[t]
    \centering
    \includegraphics[width=0.8\linewidth]{images/feature_importance_all_features_memory_rdbms_vs_qiskit.jpg}
    \caption{Importance of top 10 InferQ numeric features for SVM and Random Forest trained for optimizing memory.}
    \label{fig:memoptallrdbmsimp}
       % \vspace{-0.4cm}
\end{figure}

\subsubsection{All \sys features.}
We then include the full feature set from Section~\ref{ssec:features}.
For runtime (Figure~\ref{fig:timeoptallrdbmsimp}), for both the Linear SVM and XGBoost, 
features associated with state growth and interaction intensity are more important, including
\feat{shannon_entropy}, \feat{width}, and \feat{edge_count}, together with the
SQL-derived features. With XGBoost, in addition,  graph-related features are important, 
such as \feat{average_clustering_coefficient}, 
\feat{average_shortest_path_length}, and \feat{average_degree}, indicating that irregular interaction structure matters for
runtime differences.

For peak memory (Figure~\ref{fig:memoptallrdbmsimp}), the Linear SVM and Random Forest agree on
the main features: \feat{num_qubits} and \feat{shannon_entropy} dominate,
with additional contributions from \feat{width}, \feat{edge_count}, degree statistics, and
sparsity measures. The Random Forest captures non-linear thresholds in these features,
separating regimes where sparse relational execution substantially reduces peak memory. 
% \rh{XG boost is dominated by num\_qubits, however it is not a better selector than RF and thus not essential to evaluate these features despite the high accuracy}.

% Overall, the most influential features consistently describe effective interaction density and
% state growth, aligning with when relational (sparse/factorized) simulation is expected to be
% advantageous over dense statevector-style simulation.

\revisionB{\subsubsection{Ablation study}
\label{sssec:abl}
To quantify the contribution of each feature group, we further run a feature-group
ablation on the routing task. Table~\ref{tab:feature_group_ablation_main} reports  the results for XGBoost on the 7705 dataset. The full results with all models are reported in Appendix~B 
% \ref{appendix:abl} 
of \cite{inferqtech}. In Section~\ref{ssec:sqlonly}, we have observed that SQL-only features already provide a useful
indication, especially for runtime routing, showing that the shape of the generated SQL query
workload is informative. Here, we see that
the full feature set gives the best overall accuracy and F1 score.}

\begin{table}[t]
\centering
\small
\caption{\revisionB{Feature-group ablation for routing (XGBoost)}}
\label{tab:feature_group_ablation_main}
\begin{tabular}{lcccc}
\toprule
\multirow{2}{*}{Feature set} &
\multicolumn{2}{c}{Runtime} &
\multicolumn{2}{c}{Memory} \\
\cmidrule(lr){2-3}\cmidrule(lr){4-5}
 & Acc. & F1 & Acc. & F1 \\
\midrule
\textbf{Static + Graph + SQL + Dynamic} & \textbf{0.95} & \textbf{0.84} & \textbf{0.98} & \textbf{0.98} \\
Static + Graph + SQL & 0.93 & 0.78 & 0.97 & 0.96 \\
Static + Graph & 0.92 & 0.74 & 0.96 & 0.95 \\
Static & 0.92 & 0.74 & 0.95 & 0.94 \\
\bottomrule
\end{tabular}
\end{table}

\medskip
\noindent
\setlength{\fboxsep}{6pt}
\fcolorbox{black!15}{black!4}{%
  \parbox{\dimexpr\linewidth-2\fboxsep-2\fboxrule\relax}{%
\para{Take-away}
\textbf{(1)} SQL-only features are already strong predictors: the SQL query workload shape (joins, predicates, aggregates) explains much of when an RDBMS should be chosen.
\emph{This enables follow-up work on SQL-level cost modeling and query optimization tailored to quantum circuit simulation workloads.}\\[2pt]
\textbf{(2)} When using the full feature set, the top features are consistent with the underlying bottlenecks: state growth indicators such as entropies, circuit scale like qubit count, and interaction structure (graph features).
    \emph{This supports future work on data-centric models that learn when to use an RDBMS engine and how to optimize execution from \sys-generated training data.}
}}%
 \vspace{0.2cm}

% Some conclusions that are directly drawn are that the number of JOIN, AND and AGGREGATE in the SQL query are a determining factor for choosing RDBMS. The fewer of these intrinsically expensive JOIN database operations we have the better execution time the RDBMS system will have for the circuit (as num\_joins has positive importance for time optimality). While, higher number of AND and AGGREGATES are determining factors of memory and time optimality and are favored by RDBMS (SQLite), with its negative importance. It underscores how we can leverage fast and efficient operations like AGGREGATE and FILTERS (WHERE clauses) over python frameworks like Qiskit Aer. Furthermore, \textit{justifying the use of this SQL representation through showcasing how RDBMS systems can be exploited for more efficient classical simulation of quantum computation}.




% Given that this experiment shows the ability to utilize the SQL representation of a quantum circuit and selecting whether to use an RDBMS system or not, we can extend it by solving the mutliclass problem and choosing which specific Qiskit Aer method is optimal if not RDBMS.

% InferQ makes such opportunities actionable by exposing circuit structure, sparsity proxies, and interaction patterns as first-class features.
% A promising new runtime for classical simulation is RDBMS systems. Specifically,  translating quantum circuits into SQL queries can enable these RDBMS simulators to outperform alternatives like Qiskit for sparse circuits \cite{hai2025quantum}. With InferQ, we want to be able to use inexpensive features that are not complex like the dynamic metric sparsity and a data-centric approach to answer this question.    




